% ===== sec/07_methodology.tex =====
\section{Methodology: The Goal-Driven\,/\,Data-Driven Architecture}
\label{sec:methodology}

This section describes the architecture in operational terms.
The architecture is austere by design: one Goal-Driven Type-0
generator, one Data-Driven Type-3 verifier, and an explicit
recursive enumeration procedure that maps the start language to
the realised trace language.  The Type-2 and Type-1 sub-agents
appear as derivations of the binding, not as independent design
choices.

\subsection{The two principal agents}

\paragraph{Goal-Driven generator $\mathcal{G}$.}
$\mathcal{G}$ is a Type-0 agent realised concretely as a
large-language-model loop over the union of the operation
alphabets in \texttt{trace\_language\_original.csv}.  Its
configuration space is the LLM's context window plus a finite
set of bookkeeping flags; its step relation
$\delta_\mathcal{G}$ samples a continuation, parses the
continuation's tool-call markup, and emits the corresponding
operation string.  The model used at run time is
\texttt{invoke\_gemini(gemini-2.5-pro)} for the Type-0 deliberation
steps and \texttt{invoke\_gemini(gemini-2.5-flash-lite)} for the
finite-state housekeeping steps.

$\mathcal{G}$ is given exactly the following information at
start-up:
\begin{itemize}
\item the start-language alphabet
$\Sigma_0$ (the initial seed template operations);
\item the empty workspace directory \texttt{/comp\_dir/};
\item the boilerplate evaluation harness configuration.
\end{itemize}
$\mathcal{G}$ is \emph{not} given the goal label
(\textsf{Transported}), the public/private leaderboard
specification, or any prior solution.  Its trace language is, in
principle, the entirety of $\Sigma_0^*$, restricted only by what
$\mathcal{G}$'s LLM core happens to enumerate.

\paragraph{Data-Driven verifier $\mathcal{V}$.}
$\mathcal{V}$ is the deterministic pipeline state DFA described in
Section~\ref{sec:indexed} and Proposition~\ref{prop:verifier-dfa}.
It tracks the compilation, optimization, and verification lifecycle
of the pipeline graph.
$\mathcal{V}$ has no LLM component; it is a deterministic linear-time
checker.  Its decision function is formally the \emph{success criterion}:
\[
h_\mathcal{V}(\tau) \;=\;
\begin{cases}
\mathsf{accept} & \text{if the DFA ends in } \textsf{SUBMITTED} \text{ or } \textsf{PACKAGED} \\
\mathsf{reject} & \text{otherwise.}
\end{cases}
\]
A trace $\tau$ is strictly successful iff $h_\mathcal{V}(\tau) = \mathsf{accept}$.

\subsection{The binding}

The architecture's single load-bearing claim is that
\[
L(\mathcal{G}\,\Vert\,\mathcal{V}) \;=\; L(\mathcal{G})\;\cap\;
L(\mathcal{V}) \quad\in\quad \mathcal{L}_0
\]
with membership decidable in $O(n)$ on the verifier side
(Theorem~\ref{thm:closure}, Corollary~\ref{cor:budget}).  The
choice of Type-3 as $\mathcal{V}$'s class is not a free parameter:
it is the principled floor of the verifier-expressiveness trade-off
established in Section~\ref{sec:verifier-ladder}.

In implementation, the binding is realised by an outer loop:
\begin{enumerate}
\item $\mathcal{G}$ proposes a candidate trace prefix $\tau_t$.
\item $\mathcal{V}$ runs its DFA on $\tau_t$, advancing its state in
linear time.
\item If $\mathcal{V}$ transitions to \textsf{REJECTED} or \textsf{ERROR},
the rejected suffix is discarded
and $\mathcal{G}$ is advised (by feeding back the failure state
as a percept) to enumerate further.
\item If $\mathcal{V}$ accepts and the accepted trace contains
a \textsf{SUBMIT} step, the loop terminates and
the graph is uploaded.
\end{enumerate}
The loop terminates because the laptop's tape is finite (see
Section~\ref{sec:discussion}), or because the verifier accepts a
trace ending in \textsf{yield\_final\_result}, whichever comes
first.

\subsection{Derivation of the Type-2 and Type-1 sub-agents}

The Type-2 sub-agents \texttt{analyzer},
\texttt{miner} and the Type-1 sub-agent
\texttt{optimizer} were not explicit components
designed independently of the problem logic.  They are projections
of the verified trace language
$L(\mathcal{G}\,\Vert\,\mathcal{V})$ onto the orchestration
substrate.

\begin{definition}[Projection of a trace language]
\label{def:projection}
Let $\mathcal{A}$ be an agent with operation alphabet
$\Sigma_\mathcal{A}$ and trace language $L(\mathcal{A})\subseteq
\Sigma_\mathcal{A}^*$.  Let $\Sigma' \subseteq
\Sigma_\mathcal{A}$ be a sub-alphabet.  The \emph{projection}
of $L(\mathcal{A})$ onto $\Sigma'$ is the set
\[
\pi_{\Sigma'}(L(\mathcal{A})) \;=\;
\{\,\pi_{\Sigma'}(\tau) \,:\, \tau\in L(\mathcal{A})\,\}
\]
where $\pi_{\Sigma'}: \Sigma_\mathcal{A}^*\to{\Sigma'}^*$
deletes every symbol not in $\Sigma'$.
\end{definition}

\begin{theorem}[Derived sub-agents]
\label{thm:derived}
Let $\Sigma_{\textsf{analyzer}}$ be the operation alphabet of the
\texttt{analyzer} column of \texttt{experiment\_6/trace\_language.csv}.
Then
\[
L(\texttt{analyzer}) \;=\;
\pi_{\Sigma_{\textsf{analyzer}}}\bigl(\,L(\mathcal{G}\,\Vert\,\mathcal{V})\,
\bigr).
\]
The class of $L(\texttt{analyzer})$ is no greater than
$\mathcal{L}_0$ (because projection cannot raise class) and is
mapped structurally in Section~\ref{sec:chomsky} to
$\mathcal{L}_2$.  Analogous statements hold for
\texttt{miner} (Type-2),
and \texttt{optimizer} (Type-1).
\end{theorem}

\begin{proof}[Sketch]
The verified trace language is the language of all interleaved
multi-agent emissions admissible under the orchestration
relation.  Projecting onto the planner's alphabet selects exactly
the planner's emitted symbols.  The class bound is the standard
projection lemma~\cite{hopcroft2006automata}; the empirical
class assignment is by structural inspection of the column.\qed
\end{proof}

The minimality slogan of the paper is therefore literally true:
the architect designs only $\mathcal{G}$ and $\mathcal{V}$.  The
intermediate-class agents emerge from the binding via the
following explicit projection step.

\begin{figure}[ht]
\noindent\rule{\columnwidth}{0.4pt}
\textbf{Algorithm 2.} \textsc{ProjectSubAgent}$(\tau, \Sigma')$
\par\noindent\rule{\columnwidth}{0.4pt}
\begin{algorithmic}[1]
\REQUIRE Verified trace $\tau$, target alphabet $\Sigma'$ for the sub-agent.
\ENSURE Projected sub-trace $\pi_{\Sigma'}(\tau)$.
\STATE $\tau_{sub} \gets \varepsilon$;
\FOR{each $w_i \in \tau$}
  \IF{$w_i \in \Sigma'$}
    \STATE $\tau_{sub} \gets \tau_{sub} \cdot w_i$;
  \ENDIF
\ENDFOR
\RETURN $\tau_{sub}$;
\end{algorithmic}
\rule{\columnwidth}{0.4pt}
\caption{Sub-agent Projection from the verified trace.}
\label{alg:projection}
\end{figure}

\subsection{The recursive-enumeration procedure}

Algorithm~\ref{alg:enum} is the procedure that maps the start
language to the realised trace language.

\begin{figure}[ht]
\noindent\rule{\columnwidth}{0.4pt}
\textbf{Algorithm 1.} \textsc{RecursivelyEnumerate}$(\mathcal{G},
\mathcal{V}, x)$
\par\noindent\rule{\columnwidth}{0.4pt}
\begin{algorithmic}[1]
\REQUIRE Goal-Driven $\mathcal{G}$, Data-Driven verifier
$\mathcal{V}$, input $x$.
\ENSURE Verified trace $\tau\in L(\mathcal{G}\,\Vert\,\mathcal{V})$,
or \textsf{halt-on-tape-bound}.
\STATE $q\gets q_0^\mathcal{G}$;\quad $\tau\gets\varepsilon$;\quad
$s\gets q_0^\mathcal{V}$.
\LOOP
  \STATE $(q', w) \gets \delta_\mathcal{G}(q, x)$.
  \STATE $s' \gets \delta_\mathcal{V}^*(s, w)$ \,\,\,\,\,\#
       advance the pipeline DFA on $w$.
  \IF{$s'$ is a reject/error state}
    \STATE feed the failure witness back to $\mathcal{G}$
    as a percept.
  \ELSE
    \STATE $\tau \gets \tau\cdot w$;\quad $q\gets q'$;\quad $s\gets s'$.
    \IF{$q\in F_\mathcal{G}$ and $s\in F_\mathcal{V}$}
      \RETURN $\tau$.
    \ENDIF
  \ENDIF
  \IF{laptop tape exhausted}
    \RETURN \textsf{halt-on-tape-bound}.
  \ENDIF
\ENDLOOP
\end{algorithmic}
\rule{\columnwidth}{0.4pt}
\caption{Recursive enumeration of the start language.}
\label{alg:enum}
\end{figure}

The structure of Algorithm~\ref{alg:enum} is exactly the structure
of a Type-0 step (line~3) intersected with a Type-3 advance
(line~4).  By Theorem~\ref{thm:closure}, the loop's output
language $L(\mathcal{G}\,\Vert\,\mathcal{V})$ is at most
$\mathcal{L}_0$ and at least $\mathcal{L}_3$ (the verifier's
class).  By the laptop's finite tape, the loop terminates in
finite time.

\subsection{What the algorithm produces}

When Algorithm~\ref{alg:enum} is run on the
\textsc{NeuroGolf ONNX} input, the operation alphabet
$\Sigma_0$ extends to $\Sigma_1 = \Sigma_0\cup\{\text{symbols
emitted along the verified path}\}$.  The additional columns of
the resulting CSV are the emergent agents:
strictly new sub-agents whose alphabets contain symbols absent
from $\Sigma_0$ (e.g.,
\texttt{FP16\_SURGERY}, \texttt{BLEND\_BUNDLE}, \texttt{VERIFY\_TRAIN}).
Their emergence is forced by the verifier rejecting any trace that
fails to optimize or correctly shape the ONNX graph, not by the architect having pre-designed
them.

The empirical CSV is therefore a fixed point of the binding:
\[
\Sigma_1 \;=\;
\Sigma_0 \cup \mathrm{Range}\bigl(\delta_\mathcal{G}\bigr)
\big|_{\mathcal{V}\text{-accepting}}
\]
and the verified trace $\tau^*\in L(\mathcal{G}\,\Vert\,\mathcal{V})$
is exactly the column-wise concatenation of the recorded entries
in \texttt{experiment\_6/trace\_language.csv}.

\subsection{Hoare-triple obligation per step}

Each step of Algorithm~\ref{alg:enum} carries a Hoare-triple
obligation in the sense of
Hoare~\cite{hoare1969axiomatic} and
Dijkstra~\cite{dijkstra1975guarded}:
\[
\{\,P_t\,\}\quad w_t\quad\{\,P_{t+1}\,\}
\]
where $P_t$ is the DFA state before the $t$-th operation symbol
$w_t$, and $P_{t+1}$ is the DFA state after the transition
$\delta_\mathcal{V}(P_t, w_t)$.  For concreteness: the triple
\[
\{\,\textsf{ANALYZED}\,\}\;\textsf{BUILD\_ONNX}\;\{\,\textsf{BUILT}\,\}
\]
expresses that an \textsf{ANALYZED} pipeline may be built;
violating the precondition (e.g., attempting \textsf{BUILD\_ONNX}
from \textsf{INIT}) drives the DFA to \textsf{ERROR}.  The DFA
transition table is the \emph{mechanical realisation} of these
obligations.  We note explicitly that this is a structural analogy:
the DFA does not execute a typed pre/post-condition verifier; it
executes a finite-state membership check that \emph{plays the role}
of the Hoare obligation at the pipeline level.

\subsection{Temporal-logic statement of correctness}

The intended correctness properties of the binding are captured as
LTL~\cite{pnueli1977temporal} formulas below.  These are
\emph{design-level specifications}, not machine-checked proofs;
no Kripke structure or model-checker was run.
\begin{align*}
\textsf{Safety:}\quad &
G\bigl(\,\mathsf{emit}\,\to\,\mathsf{verifier\text{-}tracking}\bigr) \\
\textsf{Liveness:}\quad &
G\bigl(\,\mathsf{plan}\,\to\,F\,(\mathsf{verified}\,\vee\,
\mathsf{halt\text{-}on\text{-}tape\text{-}bound})\bigr).
\end{align*}
Safety states that every emitted symbol advances the verifier's
state.  Liveness states that every plan eventually either succeeds
verification or hits the tape bound.  Both are satisfied by
construction when Algorithm~\ref{alg:enum} is applied to a
finite tape: the DFA either advances or rejects; the loop either
terminates on a verified trace or on the tape bound.  Formal
model-checking on a finite Kripke abstraction of the DFA is
left as future work~\cite{clarke1981ctl,clarke1999modelchecking}.

This concludes the formal description of the architecture.  We
turn next to the empirical instantiation.
